% Applied & Computational Mathematics: Key Concepts
% Extracted from introduction.tex and preparation_guide.tex

\documentclass{amsart}
\usepackage{amsmath, amssymb, amsthm}
\usepackage{enumitem}

\title{Applied \& Computational Mathematics: Key Concepts}
\author{S.-T. Yau College Student Mathematics Contests}
\date{Last updated: May 2025}

\begin{document}

\maketitle

\section{Calculus (syllabus)}

\begin{itemize}
  \item Functions, limits, continuity
  \item Derivatives: tangent lines, optimization, related rates
  \item Integrals: indefinite, definite, improper; FTC; techniques (substitution, parts, partial fractions, trig substitution)
  \item Sequences and series; convergence tests; power series; Taylor series
  \item Multivariable: partial derivatives, gradient, divergence, curl, Jacobian
  \item Multiple integrals; line and surface integrals
  \item Green, Gauss, Stokes theorems
  \item ODEs: separable, linear, exact, integrating factor
  \item Higher-order linear ODEs with constant coefficients
  \item Systems of first-order ODEs; phase portraits
  \item PDEs: classification (elliptic, parabolic, hyperbolic); canonical equations (heat, wave, Laplace)
\end{itemize}

\section{Linear Algebra (syllabus)}

\begin{itemize}
  \item Determinants; cofactor expansion
  \item Matrix operations; inverse matrices; Sherman--Morrison rank-one update
  \item Gaussian elimination; LU; RREF
  \item Vector spaces: subspaces, basis, dimension, coordinates
  \item Linear independence; rank--nullity
  \item Eigenvalues and eigenvectors; characteristic polynomial; Cayley--Hamilton
  \item Spectral theorem for symmetric matrices
  \item Inner product; Gram--Schmidt; QR decomposition; least squares
  \item Quadratic forms; Gershgorin disc theorem
  \item SVD; low-rank approximation via SVD; Eckart--Young
  \item Condition numbers of augmented systems for least squares
  \item Pseudoinverse bounds $\|x\| \leq \|\mathbf{b}\|/\sigma$
  \item Diagonal dominance conditions for matrix invertibility
  \item Properly Hessenberg matrices and Krylov subspaces
\end{itemize}

\section{Probability (syllabus)}

\begin{itemize}
  \item Sample spaces, events, probability axioms; conditional probability; Bayes
  \item Random variables: discrete and continuous distributions
  \item Joint distributions; marginal and conditional
  \item Expectation, variance, covariance, correlation; MGF and CF
  \item LLN; CLT; conditional expectation
\end{itemize}

\section{Operations Research (syllabus)}

\begin{itemize}
  \item LP: formulation, geometry, simplex, two-phase, big-M, duality; sensitivity analysis
  \item Integer LP; branch-and-bound
  \item Transportation and assignment; Hungarian algorithm
  \item Network flows: Dijkstra, Floyd--Warshall, Kruskal, Prim, Ford--Fulkerson
  \item CPM/PERT; decision analysis: decision trees, expected value, risk
  \item Game theory: zero-sum, Nash equilibrium, mixed strategies
  \item Dynamic programming for investment optimization
  \item Resource allocation problem via dynamic programming
\end{itemize}

\section{Finite Difference Methods for PDEs}

\begin{itemize}
  \item Consistency, stability, convergence; Lax equivalence theorem
  \item Von Neumann stability analysis; CFL condition
  \item Explicit schemes: FTCS for heat equation; stability constraints
  \item Upwind scheme; Beam--Warming; Lax--Wendroff; leapfrog scheme
  \item Implicit schemes: BTCS; Crank--Nicolson ($\theta$-scheme); unconditional stability
  \item Richardson extrapolation
  \item Finite difference for wave equation and Laplace/Poisson
  \item Lax-Friedrichs scheme and its instability without diffusion
  \item DuFort-Frankel scheme
  \item Three-point scheme for $u_t + au_x = f$
  \item Richardson's scheme for heat equation (unconditionally unstable)
  \item Implicit leapfrog scheme (order $(2,4)$)
  \item Finite difference approximations for $-\frac{d}{dx}(a(x)u') = f$
  \item $\theta$-scheme for $u_t + u = u_{xx}$ ($\ell_\infty$ and $\ell_2$ stability)
  \item WKB expansion for Schr\"odinger equation
  \item Well-posedness of $u_t = u_{xxx}$ via Fourier transform
  \item Diffusion equation stability: FE ($k \leq h^2/(2\mu)$), CN (unconditional)
  \item $L^2$ stability of semi-discrete upwind scheme for conservation laws
  \item $L^2$ stability for finite differences applied to $-u'' + \alpha u = f$
  \item Stability of perturbations ($B + \Delta t C$ remains stable)
\end{itemize}

\section{Finite Element Methods (FEM)}

\begin{itemize}
  \item Weak (variational) formulation of BVPs; Ritz--Galerkin
  \item CG and DG methods; piecewise polynomial spaces ($P_1$, $P_2$)
  \item A priori error estimates; convergence rates in $H^1$ and $L^2$
  \item Superconvergence at Gaussian points
  \item FEM for elliptic BVPs: Poisson equation
  \item Error estimates for $L^2$ projection into piecewise polynomial spaces
  \item Superconvergence of FEM for 1D Poisson (exactness at nodes)
  \item Immersed interface method for elliptic equations with discontinuous coefficients
  \item $L^2$ stability of discontinuous Galerkin method for conservation laws
  \item Beam equation $u^{(4)} = f$ with Cantilever BCs
\end{itemize}

\section{Numerical Linear Algebra}

\begin{itemize}
  \item SVD; low-rank approximation via SVD; Eckart--Young
  \item QR: Gram--Schmidt (classical and modified), Householder, Givens
  \item Vandermonde least-squares approximation
  \item Eigenvalue algorithms: power method; inverse iteration; shifted power method; QR algorithm
  \item Arnoldi process; Krylov subspaces; Lanczos
  \item CG; GMRES; Richardson; SOR; Jacobi; Gauss--Seidel
  \item Spectral radius of tridiagonal (Toeplitz) matrices; Gershgorin discs
  \item Spectral radius of tridiagonal matrix of all ones via Chebyshev polynomials
  \item Power method convergence rate $|\lambda_2/\lambda_1|$
  \item Inverse iteration and eigenvector sensitivity analysis
  \item Subspace iteration and Rayleigh-Ritz approximation
  \item Constructing orthonormal bases from Lagrange polynomials at Gaussian nodes
  \item Nuclear norm minimization for low-rank approximation
  \item PCA/CCA for finding optimal orthonormal vectors
  \item Optimal gradient method for linear systems
  \item 2D fixed-point iteration convergence (Gauss-Seidel-type modification)
  \item Stationary iterative methods ($x^{k+1} = Ax^k + b$); spectral radius conditions
  \item eigenvalues of augmented system matrix $\begin{bmatrix} I & X \\ X^T & 0 \end{bmatrix}$
\end{itemize}

\section{Numerical Methods for ODEs}

\begin{itemize}
  \item Forward Euler, backward Euler, improved Euler, RK4
  \item Consistency, order, convergence; absolute stability regions
  \item BDF; BDF2 zero-stability and order; SSP-RK; IMEX RK
  \item Multistep: Adams--Bashforth, Adams--Moulton
  \item Stiff ODEs; A-stability, L-stability
  \item Banach fixed-point for ODE convergence
  \item Secant method (superlinear convergence); Steffensen acceleration; bisection; Newton's method
  \item Steffensen's acceleration method for fixed point iteration (quadratic convergence)
  \item General one-step Runge-Kutta method; order conditions $\alpha+\beta=1$, $\beta\gamma=1/2$
  \item Semi-implicit Runge-Kutta methods; A-stability for $\beta > 1/2$
  \item Strong stability preserving (SSP) Runge-Kutta methods; optimal coefficients
  \item Velocity-Verlet method as symplectic map
  \item St\"ormer's method (order 2) for wave equation; CFL stability condition $\Delta t / \Delta x \leq 1$
  \item Implicit scheme for $u_t = H(u)_{xx}$
  \item Multistep methods for $u_t + au_x = f$ (unconditional stability)
  \item IMEX time discretization for convection-diffusion ($L^2$ stability under $\Delta t \leq 2b/a^2$)
  \item Third-order Lax-Wendroff-type schemes for $u_t + au_x = 0$
  \item Multipoint iteration for root-finding (Steffensen-type method); cubic convergence
  \item Local error bound for $y' = -\lambda y$
  \item Stability regions for explicit Euler, implicit Euler, trapezoidal rule
\end{itemize}

\section{Optimization and Convex Analysis}

\begin{itemize}
  \item KKT conditions; complementary slackness; dual feasibility
  \item Lagrangian duality; weak and strong duality; Slater's condition
  \item Lasso regression; $L^1$ regularization; compressed sensing
  \item Convex sets and functions; co-coercivity; proximal operators; FBP; Douglas--Rachford
  \item Subgradient and subdifferential calculus; Frank--Wolfe basics
  \item Descent Lemma and co-coercivity inequalities for convex functions
  \item Strong convexity; equivalence of three definitions
  \item Metric projection onto convex sets; variational inequality; nonexpansivity
  \item Gradient descent for convex functions
  \item Backward Euler for convex minimization with strongly convex functions
  \item Uncertainty principle $\|x\|_0 \|\hat{x}\|_0 \geq N$
\end{itemize}

\section{Calculus of Variations and Euler--Lagrange}

\begin{itemize}
  \item Euler--Lagrange equation; necessary and sufficient conditions; natural boundary conditions
  \item Isoperimetric problems; isoperimetric inequality
  \item Legendre condition; Jacobi conjugate points
  \item Noether's theorem; direct method; Ritz method
  \item Similarity solutions for porous medium equation $h_t = \Delta(h^2)$
  \item Finite time extinction and hyper-contractivity for ODEs ($y' \leq \alpha - \beta y^a$)
  \item Scaling laws and slowly varying functions
\end{itemize}

\section{Orthogonal Polynomials and Spectral Methods}

\begin{itemize}
  \item Legendre polynomials: orthogonality, Rodrigues formula, recurrence
  \item Associated Legendre functions
  \item Chebyshev ($T_n$, $U_n$): orthogonality with weight $(1-x^2)^{-1/2}$; minimax property
  \item Chebyshev polynomials of second kind $U_n(x) = \sin((n+1)\theta)/\sin\theta$
  \item Hermite polynomials: orthogonality, recurrence $H_{n+1} = xH_n - nH_{n-1}$, eigenfunction property $xu' - u'' = \lambda u$
  \item Gaussian quadrature: Gauss--Legendre, Gauss--Chebyshev; error estimates
  \item Clenshaw--Curtis; Lobatto; Radau
  \item Chebyshev interpolation; spectral collocation
  \item Moving least squares (MLS); error analysis
  \item Fourier spectral methods; spectral accuracy
  \item Crank--Nicolson spectral for diffusion
  \item Envelope of $T_{n+1}(x) - T_{n-1}(x)$ forming an ellipse
  \item Chebyshev polynomial properties: $T_n(x) = \cos(n\arccos x)$, extrema at $\cos(k\pi/n)$, leading coefficient $2^{n-1}$
  \item Orthogonal polynomials for weight $\log(1/x)$; moment calculations; recurrence relations; Jacobi matrix eigenvalues
  \item Two-point Gaussian quadrature
  \item Interpolation error bounds; cubic spline properties
  \item Polynomial approximation error propagation across subintervals; Markov-type inequalities
  \item Ballot theorem and lattice path enumeration $\frac{1}{r+s}\binom{r+s}{s}$
\end{itemize}

\section{Dynamical Systems and Biological Modeling}

\begin{itemize}
  \item Phase plane; fixed points; linearization via Jacobian; stability; Lyapunov stability
  \item Bifurcations: saddle-node, transcritical, pitchfork, Hopf
  \item FitzHugh--Nagumo neuronal model; excitability; quiescent vs. tonic spiking states
  \item Michaelis--Menten enzyme kinetics; steady-state approximation; Lineweaver--Burk plot; singular perturbation
  \item WKB for Schr\"odinger equations
  \item Hamilton--Jacobi equations; shock formation; Rankine--Hugoniot conditions
  \item Fokker-Planck, Witten Laplacian, and Schr\"odinger operators; eigenvalue equivalence
  \item Energy estimates for backward Euler discretization of heat equation
  \item Energy conservation for wave guide problem $u_t + u_x = 0$, $v_t - v_x = 0$
  \item Weak solutions of wave equation (distributional sense)
  \item Nondimensionalization of Stokes drag
  \item Scaling analysis of projectile motion under various limits
  \item Mark Kac's question ``Can you hear the shape of a drum?'' for 1D oscillators
  \item Eigenvalues of hypercube $C_n$ via Walsh functions
\end{itemize}

\section{Graph Theory}

\begin{itemize}
  \item Eulerian and Hamiltonian graphs; Dirac's theorem ($n\geq3$, $\deg(v)\geq n/2$); Ore's theorem
  \item Bondy--Chvatal closure
  \item Hall's marriage theorem; Tutte's theorem
  \item Planar graphs; Kuratowski ($K_5$, $K_{3,3}$)
  \item Graph coloring; chromatic polynomial
  \item Hamilton-connected graphs; degree sum $\geq n+1$ implies Hamiltonian path
  \item $(p,q)$-condition for matchings in bipartite graphs
  \item Structural balance theory in social networks
\end{itemize}

\section{Number Theory}

\begin{itemize}
  \item Divisibility; Euclidean algorithm; prime numbers
  \item Congruences; Chinese remainder theorem
  \item Fermat's little theorem; Euler's theorem; $\varphi(n)$
  \item RSA encryption; primitive roots; discrete logarithms
  \item Quadratic residues; Legendre and Jacobi symbols
  \item Eisenstein's criterion; Eisenstein polynomials and multivariate generalizations
  \item Efficient gcd computation for polynomials via modular exponentiation
  \item Efficient RSA factorization for twin primes
  \item Sieving methods for multiplicative dependence; bounds on exponent vectors
  \item Siegel's lemma on small integer solutions to homogeneous linear systems
  \item Generating functions; Catalan number connections
  \item Half-angle formula with Catalan numbers
  \item Complete monotonicity of moment sequences (Bernstein's theorem)
  \item Minkowski's existence theorem for convex polyhedra
  \item Branching processes and extinction probabilities
\end{itemize}

\section{Learning Theory and VC Dimension}

\begin{itemize}
  \item PAC learning framework; sample complexity
  \item VC dimension: definition, growth function, Sauer's lemma
  \item VC dimension of interval classifiers in $\mathbb{R}$ equals 2
  \item VC dimension of linear classifiers in $\mathbb{R}^d$ equals $d+1$
  \item Rademacher complexity basics
  \item DFT uncertainty
\end{itemize}

\end{document}
